\documentclass[11pt,a4paper]{article}

\usepackage[utf8]{inputenc}
\usepackage{amsmath, amssymb, amsthm}
\usepackage{graphicx}
\usepackage{pgfplots}
\pgfplotsset{compat=1.18}
\usepackage{hyperref}
\usepackage{listings}
\lstset{basicstyle=\ttfamily\small, breaklines=true}
\usepackage{geometry}
\geometry{margin=1in}

% Theorem environments
\newtheorem{theorem}{Theorem}[section]
\newtheorem{definition}{Definition}[section]
\newtheorem{lemma}{Lemma}[section]
\newtheorem{proposition}[theorem]{Proposition}
\newtheorem{example}{Example}[section]

\title{\textbf{The Cartography of Paradoxes: \\ Unifying Probabilistic Epistemic Logic and Non-Bivalent Validity}}
\author{Julian L. Voigt\thanks{ORCID: \href{https://orcid.org/0009-0007-2147-3340}{0009-0007-2147-3340} \\ Final Draft}}
\date{\today}

\begin{document}

\maketitle

\begin{abstract}
Probabilistic Epistemic Logic (PEL) traditionally operates under the assumption of classical bivalent semantics, which renders it vulnerable to epistemic paradoxes such as the Preface Paradox and logical omniscience issues. Simultaneously, non-bivalent logics (e.g., Belnap-Dunn 4-valued logic) successfully resolve semantic paradoxes like the Liar but lack a robust probabilistic framework for doxastic attitudes. This paper introduces 4-Valued Probabilistic Epistemic Logic (4-PEL), a novel framework that models non-bivalent probability over a multi-agent Kripke-style structure $(W, \mu, v)$ with a 4-valued valuation function $v(w, \phi) \in \{T, B, N, F\}$. By defining an agent-specific Lockean threshold belief operator over this space, I formulate a strict mathematical boundary for the emergence of epistemic gluts (paradoxes). I provide Lean 4-verified formal proofs of the central probabilistic bounds demonstrating that epistemic gluts require a strict lower bound of ontological gluts, alongside general theorems ensuring the preservation of classical logic in paradox-free states and the prevention of epistemic explosion. I establish the computational tractability of the framework, proving a polynomial-time bound for model checking. Furthermore, I apply non-bivalent validity standards (Strict-Tolerant, Tolerant-Tolerant, Strict-Strict) to demonstrate how 4-PEL successfully models and classifies five classical epistemic paradoxes (the Liar, Lottery, Preface, Moore's Paradox, and G\"odel-inspired epistemic fixed points) while blocking epistemic explosion.
\end{abstract}

\section{Introduction}

The intersection of probability and formal epistemology has long been a fertile ground for both deep insights and intractable paradoxes. At the heart of this intersection lies the Lockean Thesis, which posits that a rational agent believes a proposition if and only if their subjective probability in that proposition exceeds a certain high, agent-specific threshold $c_i \in (0.5, 1]$. While elegantly simple, implementing this thesis within standard Probabilistic Epistemic Logic (PEL) inevitably leads to structural friction when confronting the bounds of classical logic. Specifically, classical PEL breaks down entirely when encountering self-referential paradoxes like the Liar, or aggregation issues like the Lottery Paradox.

The central result of this paper is that subjective, epistemic contradictions (epistemic gluts) are mathematically impossible unless the objective probability of an ontological contradiction (an ontological glut, $P_B$) meets a strict lower bound. Specifically, I formally derive the \textbf{Glut Boundary Theorem}, which proves that an epistemic glut for agent $i$ requires an ontological contradiction level of $P_B \ge 2c_i - 1$. 

This theorem matters because it transforms paradoxes from binary, explosive logic failures into a continuous, navigable terrain where the degree of subjective contradiction is strictly bounded by the underlying uncertainty in the objective state space. By establishing this mathematical limit, the logic explicitly maps the threshold of rational belief onto the structure of paraconsistent probability.

To achieve this, I propose the \textbf{4-Valued Probabilistic Epistemic Logic (4-PEL)}. By discarding the bivalent assumption of the underlying state space, 4-PEL allows the probability measure $\mu$ to distribute mass directly over the four Belnap-Dunn states: Strict Truth, Strict Falsity, Ontological Gaps, and Ontological Gluts. Consequently, belief is no longer a simple summation over true worlds, but an aggregation of positive evidential support.

The primary contributions of this paper are fourfold:
\begin{enumerate}
    \item I mathematically derive and formally verify (using the Lean 4 interactive theorem prover) the ``Cartography of Paradoxes'': the Glut Boundary Theorem bounding epistemic gluts by $P_B \ge 2c_i - 1$.
    \item I formalize the 4-PEL framework, integrating world-relative probability measures over 4-valued spaces with an agent-specific Lockean belief operator, and introduce dynamics for belief conditionalization.
    \item I demonstrate that mapping non-bivalent validity standards (ST, LP, SS) onto this probabilistic space successfully models and classifies five paradoxical case studies (Liar, Lottery, Preface, Moore's Paradox, and G\"odel-inspired epistemic fixed points) while blocking explosive consequences in the relevant paraconsistent settings.
    \item I prove the ``Disagreement Theorem,'' establishing how multi-agent structures natively support rational disagreement and higher-order perspective modeling. This lays a theoretical foundation for applications in social epistemology, rational disagreement, and AI ensemble alignment.
\end{enumerate}

\section{Related Work}
\label{sec:relatedwork}

\subsection{Classical PEL}
Traditional Probabilistic Epistemic Logic (PEL) evaluates propositions over a set of possible worlds $W$, where each world $w \in W$ assigns a classical, bivalent truth value (either True or False) to every proposition. A probability measure $\mu$ is defined over subsets of $W$, representing the agent's degrees of belief. The cornerstone of PEL is the \textit{Lockean Thesis}, which establishes a quantitative threshold $c$ for qualitative belief. Pan and Guo \cite{pan2024probabilistic} extensively explore this framework. In the standard static paradigm, an agent believes a proposition $\phi$ (denoted $B(\phi)$) if and only if the probability of $\phi$ meets or exceeds $c$:
\begin{equation}
    B(\phi) \iff \mu(\{w \in W \mid w \models \phi\}) \ge c
\end{equation}
While intuitive, this strict bivalence restricts the logic's ability to model semantic contradictions and ambiguity. In contexts where probabilities are non-additive or imprecise (such as Dempster-Shafer theory \cite{shafer1976}), strict bivalence is often relaxed, but these formalisms typically do not natively incorporate self-referential paradoxes. When an agent encounters the Liar Paradox (``This sentence is false''), classical PEL fails because there exists no bivalent world where the sentence can be evaluated consistently, leaving the probability measure undefined.

\subsection{FDE and Belnap-Dunn}
To address such paradoxes, structural proof theory has popularized the Belnap-Dunn 4-valued logic \cite{belnap1977}. Instead of a single bivalent assignment, propositions are evaluated using a bi-dimensional boolean pair $\langle v^+, v^- \rangle$, representing independent positive and negative support. This generates four distinct ontological states:
\begin{itemize}
    \item \textbf{T (Strict Truth):} $\langle 1, 0 \rangle$
    \item \textbf{F (Strict Falsity):} $\langle 0, 1 \rangle$
    \item \textbf{N (Ontological Gap):} $\langle 0, 0 \rangle$ (Neither true nor false)
    \item \textbf{B (Ontological Glut):} $\langle 1, 1 \rangle$ (Both true and false)
\end{itemize}

\subsection{Non-Bivalent Validity}
By admitting gluts and gaps, non-bivalent logics naturally accommodate paradoxical statements. However, standard consequence relations in classical logic break down when dealing with gluts (e.g., Modus Ponens fails if the premise is a glut). As detailed by Barrio, Fiore, and Pailos \cite{barrio2025non}, non-bivalent validity frameworks provide robust metalogical tools to preserve valid reasoning:
\begin{itemize}
    \item \textbf{Strict-to-Strict (SS):} If premises are strictly true (T), the conclusion must be strictly true (T).
    \item \textbf{Strict-to-Tolerant (ST):} If premises are strictly true (T), the conclusion must be tolerant (T or B). This successfully preserves classical tautologies while preventing explosion.
    \item \textbf{Tolerant-to-Tolerant (LP):} If premises are tolerant (T or B), the conclusion must be tolerant (T or B).
\end{itemize}

\subsection{Dynamic Epistemic Logic}
To model how agents revise their beliefs over time, Dynamic Epistemic Logic (DEL) provides operators that update the underlying Kripke models \cite{vanbenthem2011}. Dynamic approaches to probabilistic belief revision have been extensively tackled by Baltag and Smets \cite{baltag2008probabilistic}, who provide formal mechanisms for conditionalizing probability measures upon receiving new information. However, standard DEL frameworks generally assume the underlying evidence is bivalent.

\subsection{Paraconsistent Probability}
While approaches to Paraconsistent Probability \cite{priest1979} offer valuable insights into distributing probability over contradictory events, and frameworks for judgment aggregation \cite{list2002} tackle conflicting multi-agent votes, they frequently lack a unified Kripke-style possible worlds semantics tying dynamic threshold belief to 4-valued validities.

\subsection{Synthesis}
While these distinct frameworks successfully address specific facets of logic, none of them natively interpret the overlap as ontological glut mass inside a threshold-based epistemic semantics. Table \ref{tab:frameworks} illustrates the comparative capabilities of existing approaches. Crucially, the derivation of the boundary theorem depends entirely on the mathematical overlap between positive and negative support within a unified probability space. Classical PEL lacks the contradiction machinery, while categorical non-bivalent logics lack the probabilistic continuum. My objective is to bridge this gap by embedding these non-bivalent validity standards into a fully probabilistic Kripke-style framework, equipping it with dynamic capabilities.

\begin{table}[ht]
\centering
\resizebox{\textwidth}{!}{%
\begin{tabular}{lccc}
\hline
\textbf{Framework} & \textbf{Handles Probability} & \textbf{Handles Contradiction} & \textbf{Interprets Glut Boundary} \\
\hline
Classical PEL & \checkmark & $\times$ & $\times$ \\
FDE & $\times$ & \checkmark & $\times$ \\
LP/ST validity & $\times$ & \checkmark & $\times$ \\
Paraconsistent Probability & Partial & Partial & Unknown \\
\textbf{4-PEL} & \textbf{\checkmark} & \textbf{\checkmark} & \textbf{\checkmark} \\
\hline
\end{tabular}%
}
\caption{Comparative analysis of logical frameworks.}
\label{tab:frameworks}
\end{table}

\section{Formal Syntax and Semantics}
\label{sec:framework}

To rigorize the integration of threshold belief and paraconsistent reasoning, I define the formal syntax and semantics of 4-PEL.

\subsection{Syntax}
Let $\Phi$ be a finite set of atomic propositions and $Ag = \{1, \dots, n\}$ be a finite set of agents. The language $\mathcal{L}_{4PEL}$ is generated by the following Backus-Naur Form (BNF) grammar:
\begin{equation*}
    \phi ::= p \mid \top \mid \bot \mid \neg \phi \mid \phi \land \psi \mid B_i(\phi)
\end{equation*}
where $p \in \Phi$ and $i \in Ag$. Disjunction and material implication are defined as standard syntactic macros: $\phi \lor \psi := \neg(\neg \phi \land \neg \psi)$ and $\phi \to \psi := \neg \phi \lor \psi$. The operator $B_i(\phi)$ is read as ``agent $i$ believes $\phi$''.

\subsection{Model-Theoretic Semantics}
To break free from bivalent constraints, 4-PEL replaces the classical set of strictly bivalent worlds with a probability space defined over a set of non-bivalent states. I define a \textbf{Multi-Agent 4-PEL Model} as a Kripke-style structure equipped with a 4-valued valuation and world-relative probability measures for each agent.

\begin{definition}[4-PEL Model]
A model $\mathcal{M}$ is a tuple $\langle W, Ag, \{R_i\}_{i \in Ag}, \{\mu_{i,w}\}_{i \in Ag, w \in W}, v, c \rangle$ where:
\begin{itemize}
    \item $W$ is a non-empty, finite set of possible worlds (or states).
    \item $Ag = \{1, \dots, n\}$ is a finite set of agents.
    \item $R_i \subseteq W \times W$ is an accessibility relation for each agent $i$. Let $R_i(w) = \{w' \in W \mid w R_i w'\}$ denote the set of worlds agent $i$ considers possible from $w$.
    \item $\mu_{i,w}: \mathcal{P}(R_i(w)) \to [0, 1]$ is a local probability measure for agent $i$ at world $w$, defined strictly over the reachable worlds, such that $\mu_{i,w}(R_i(w)) = 1$.
    \item $v: W \times \Phi \to \{T, B, N, F\}$ is a 4-valued valuation function assigning a Belnap-Dunn truth value to each atomic proposition at each world.
    \item $c: Ag \to (0.5, 1]$ is an agent-specific threshold function mapping each agent to their Lockean threshold $c_i$.
\end{itemize}
\end{definition}

I extend the valuation $v$ recursively to all formulas in $\mathcal{L}_{4PEL}$. Let $\mathcal{M}, w \models^+ \phi$ denote that $\phi$ has positive support at $w$, and $\mathcal{M}, w \models^- \phi$ denote negative support. The semantic rules are:
\begin{itemize}
    \item $\mathcal{M}, w \models^+ \top \text{ and } \mathcal{M}, w \nvDash^- \top$
    \item $\mathcal{M}, w \nvDash^+ \bot \text{ and } \mathcal{M}, w \models^- \bot$
    \item $\mathcal{M}, w \models^+ p \iff v^+(w, p) = 1$
    \item $\mathcal{M}, w \models^- p \iff v^-(w, p) = 1$
    \item $\mathcal{M}, w \models^+ \neg\phi \iff \mathcal{M}, w \models^- \phi$
    \item $\mathcal{M}, w \models^- \neg\phi \iff \mathcal{M}, w \models^+ \phi$
    \item $\mathcal{M}, w \models^+ \phi \land \psi \iff \mathcal{M}, w \models^+ \phi \text{ and } \mathcal{M}, w \models^+ \psi$
    \item $\mathcal{M}, w \models^- \phi \land \psi \iff \mathcal{M}, w \models^- \phi \text{ or } \mathcal{M}, w \models^- \psi$
\end{itemize}

A proposition $\phi$ thus corresponds to non-bivalent extensions in $W$. The positive extension is $\mathcal{E}^+(\phi) = \{w \in W \mid \mathcal{M}, w \models^+ \phi\}$, and the negative extension is $\mathcal{E}^-(\phi) = \{w \in W \mid \mathcal{M}, w \models^- \phi\}$. Note that $\mathcal{E}^+(\phi)$ and $\mathcal{E}^-(\phi)$ do not form a strict partition of $W$; they overlap at ontological gluts ($B$) and leave gaps ($N$).

I redefine the epistemic Belief operator $B_i$ to aggregate probability mass over these non-bivalent extensions, restricted to the reachable worlds $R_i(w)$. An agent $i$ at world $w$ believes $\phi$ if the measure of the accessible worlds providing positive support for $\phi$ meets or exceeds the Lockean threshold $c$:

\begin{definition}[World-Relative Threshold Belief Operator]
Given a threshold $c_i \in (0.5, 1]$ for agent $i$, the belief operator $B_i(\phi)$ evaluates via the local measure $\mu_{i,w}$ over $R_i(w)$:
\begin{align}
    \mathcal{M}, w \models^+ B_i(\phi) &\iff \mu_{i,w}(\mathcal{E}^+(\phi) \cap R_i(w)) \ge c_i \\
    \mathcal{M}, w \models^- B_i(\phi) &\iff \mu_{i,w}(\mathcal{E}^-(\phi) \cap R_i(w)) \ge c_i
\end{align}
\end{definition}

Crucially, because $B_i(\phi)$ evaluates to a 4-value dynamically at every world, it can be seamlessly nested (e.g., $B_i(B_j(\phi))$) without collapsing epistemological distinctions. While Lockean threshold belief is notoriously non-factive ($B_i\phi \not\implies \phi$), this framework supports introspection properties characteristic of KD45 epistemic logic (positive introspection and meta-level threshold-state introspection). Specifically, KD45-style introspection holds under the following sufficient conditions: the accessibility relation $R_i$ is Euclidean and Transitive (ensuring $w' \in R_i(w) \implies R_i(w') = R_i(w)$) and the local measure is stable across accessible worlds ($\mu_{i,w'} = \mu_{i,w}$).

\begin{theorem}[Threshold-State Introspection under KD45 Conditions]
Under Euclidean, transitive accessibility ($R_i$) and stable local measures ($\mu_{i,w'} = \mu_{i,w}$), agents are introspective about their threshold states: if $B_i(\phi)$ has positive support at $w$, then $B_i(B_i(\phi))$ has positive support at $w$. Furthermore, if $B_i(\phi)$ fails to reach the positive threshold at $w$ (meta-level non-belief), the agent assigns probability $1$ to this failure state across their accessible worlds.
\end{theorem}
\begin{proof}
Let $w$ be the actual world. Suppose $B_i(\phi)$ has positive support at $w$, meaning $\mu_{i,w}(\mathcal{E}^+(\phi) \cap R_i(w)) \ge c$. Because $R_i$ is Euclidean and Transitive, for any $w' \in R_i(w)$, $R_i(w') = R_i(w)$. Combined with $\mu_{i,w'} = \mu_{i,w}$, the probability mass for $\phi$ is identical in all $w' \in R_i(w)$. Thus, $B_i(\phi)$ also evaluates to positive support in every $w' \in R_i(w)$. Consequently, $\mu_{i,w}(\mathcal{E}^+(B_i(\phi)) \cap R_i(w)) = 1 \ge c$, proving $B_i(B_i(\phi))$ has positive support. 

For negative introspection, we must carefully distinguish object-level negative belief from meta-level non-belief. In 4-PEL, object-level negative support $\mathcal{M}, w \models^- B_i(\phi)$ means $\mu_{i,w}(\mathcal{E}^-(\phi) \cap R_i(w)) \ge c$, which equates to believing the negation $\neg \phi$. This is distinct from simply lacking positive support ($\mu_{i,w}(\mathcal{E}^+(\phi) \cap R_i(w)) < c$). If we consider the meta-level predicate ``fails positive threshold'', and assume it holds at $w$, then by Euclidean and Transitive accessibility and measure stability, this failure is uniform across all $w' \in R_i(w)$. Therefore, the agent assigns a probability of $1 \ge c$ to the meta-level state of lacking positive support. This ensures meta-level introspection remains coherent without falsely mapping to the object-language operator $\neg B_i(\phi)$.
\end{proof}

With the connectives defined, I can observe how structural inference operates under different validity standards.

\begin{theorem}[Modus Ponens under ST and LP]
Let $\phi, \psi$ be formulas evaluated over a 4-PEL model. Modus Ponens ($\phi, \phi \to \psi \models \psi$) is \textbf{valid} under ST (Strict-Tolerant) validity, but \textbf{invalid} under LP (Tolerant-Tolerant) validity.
\end{theorem}
\begin{proof}
ST-validity requires that if the premises are strictly true ($T = \langle 1,0 \rangle$), the conclusion must be at least tolerant ($T$ or $B$, meaning positive support $v^+ = 1$). If $v(w, \phi) = T$ and $v(w, \phi \to \psi) = T$, then $v^+(\phi)=1, v^-(\phi)=0$. For implication, $\max(v^-(\phi), v^+(\psi)) = 1$. Since $v^-(\phi)=0$, we must have $v^+(\psi)=1$. Thus $\psi$ is tolerant, satisfying ST-validity.

LP-validity requires that if premises are tolerant, the conclusion is tolerant. Suppose $v(w, \phi) = B = \langle 1,1 \rangle$ and $v(w, \psi) = F = \langle 0,1 \rangle$. Both $\phi$ and $\phi \to \psi$ (since $v(w, \phi \to \psi) = \langle \max(1,0), \min(1,1) \rangle = \langle 1,1 \rangle = B$) are tolerant, but $\psi$ is strictly false ($v^+(\psi)=0$), failing LP.
\end{proof}

This structure allows me to quantitatively measure how much objective ontological inconsistency is required to sustain subjective paradoxes (Epistemic Gluts).
\section{Dynamics and Conditionalization}
\label{sec:dynamics}

A robust epistemic logic requires a mechanism for belief revision upon receiving new evidence. In classical probability, this is handled via Bayesian conditionalization. In 4-PEL, updating a world-relative measure $\mu_{i,w}$ with a 4-valued evidence proposition $E$ requires a generalized conditionalization rule. Because evidence $E$ may itself be a glut or a gap, an agent must condition over the positive extension $\mathcal{E}^+(E) \cap R_i(w)$.

\begin{definition}[Dynamic Conditionalization]
When an agent $i$ at world $w$ learns evidence $E$, the updated probability measure $\mu_{i,w}^*$ for any arbitrary subset $A \subseteq R_i(w)$ is given by the conditionalization rule:
\begin{equation}
    \mu_{i,w}^*(A) = \frac{\mu_{i,w}(A \cap \mathcal{E}^+(E) \cap R_i(w))}{\mu_{i,w}(\mathcal{E}^+(E) \cap R_i(w))}
    \label{eq:cond}
\end{equation}
provided $\mu_{i,w}(\mathcal{E}^+(E) \cap R_i(w)) > 0$. 
\end{definition}
\begin{proof}
To verify that $\mu_{i,w}^*$ is a valid probability measure over $\mathcal{P}(R_i(w))$, we must show non-negativity, finite additivity, and total mass $1$. Since $\mu_{i,w}$ is a valid measure, $\mu_{i,w}(X) \ge 0$ for all $X$. Thus, the numerator $\mu_{i,w}(A \cap \mathcal{E}^+(E) \cap R_i(w)) \ge 0$. Given the condition $\mu_{i,w}(\mathcal{E}^+(E) \cap R_i(w)) > 0$, the updated measure is non-negative. Finite additivity follows directly from the additivity of the numerator measure over disjoint sets $A_1, A_2$. 
For total mass, substitute the maximal accessible state space $A = R_i(w)$ into the rule:
$\mu_{i,w}^*(R_i(w)) = \frac{\mu_{i,w}(R_i(w) \cap \mathcal{E}^+(E) \cap R_i(w))}{\mu_{i,w}(\mathcal{E}^+(E) \cap R_i(w))} = \frac{\mu_{i,w}(\mathcal{E}^+(E) \cap R_i(w))}{\mu_{i,w}(\mathcal{E}^+(E) \cap R_i(w))} = 1$. 
Thus, dynamic conditionalization strictly preserves valid probability measures.
\end{proof}

This Bayesian rule ensures that probabilities assigned to paradoxes ($P_B$) dynamically evolve as the agent encounters new, potentially contradictory evidence. 

It is important to acknowledge alternative conditionalization strategies. For instance, one might propose conditioning positive and negative support completely independently (e.g., updating $\mathcal{E}^+(\phi)$ relative only to $\mathcal{E}^+(E)$, and $\mathcal{E}^-(\phi)$ relative only to $\mathcal{E}^-(E)$). However, independently updating these extensions risks structurally decoupling them, potentially leading to probability assignments that violate the finite additivity of the underlying worlds. By instead conditioning the entire measure $\mu_{i,w}$ exclusively on the positive support $\mathcal{E}^+(E)$, the framework treats $E$ as an assertion that the world possesses at least the positive characteristics of the evidence, mathematically preserving the coherent unity of the probability measure while keeping the dynamics tractable. Note that in this foundational formalization, the accessibility relation $R_i$ remains constant during conditionalization; incorporating full Dynamic Epistemic Logic (DEL) \cite{baltag2008probabilistic} updates, where incompatible worlds are deleted from $R_i$, is left for future work.

\section{The Core Theorem and the Cartography of Paradoxes}
\label{sec:theorem}

The most profound consequence of defining belief over a 4-valued probability space is the revelation that subjective epistemic contradictions are strictly bounded by objective ontological inconsistency. To simplify notation, let $P_v(\phi, w)$ denote the local probability mass of the worlds within the accessible state space where $\phi$ evaluates to truth value $v$, i.e., $P_v(\phi, w) = \mu_{i,w}(\{w' \in R_i(w) \mid v(w', \phi) = v\})$. I formalize this bound in the ``Cartography of Paradoxes'' Theorem.

\begin{theorem}[The Glut Boundary Theorem]
Let $\mathcal{M}$ be a 4-PEL model. For an agent $i$ at world $w$ with Lockean threshold $c_i \in (0.5, 1]$ to experience an Epistemic Glut regarding a proposition $\phi$ (i.e., $B_i(\phi)$ evaluates to both positive and negative support at $w$), it is necessary that the total local probability of the Ontological Glut state for $\phi$, denoted $P_B(\phi, w)$, is at least $2c_i - 1$. This bound is tight over the class of 4-PEL models whose accessible state space $R_i(w)$ realizes $T, F,$ and $B$ states for $\phi$: for every $x \ge 2c_i - 1$, there exists a valid local probability distribution with $P_B(\phi, w) = x$ that yields an epistemic glut.
\end{theorem}

Mathematically, this inequality is an instance of the Fr\'echet/Bonferroni lower bound $P(A \cap B) \ge P(A) + P(B) - 1$. The theoretical contribution lies not in the algebraic derivation, but in the semantic identification of the overlap term directly with the ontological glut mass ($P_B$) inside a non-bivalent epistemic semantics.

\begin{proof}
\textbf{Necessity (Epistemic Glut $\implies P_B(\phi, w) \ge 2c_i - 1$):}
By definition, $v(w, B_i(\phi))$ evaluates to a glut if $\mu_{i,w}(\mathcal{E}^+(\phi)) \ge c_i$ and $\mu_{i,w}(\mathcal{E}^-(\phi)) \ge c_i$.
This expands to $P_T(\phi, w) + P_B(\phi, w) \ge c_i$ and $P_F(\phi, w) + P_B(\phi, w) \ge c_i$.
Summing these inequalities yields $P_T(\phi, w) + P_F(\phi, w) + 2P_B(\phi, w) \ge 2c_i$.
Since probabilities must sum to 1, $P_T(\phi, w) + P_F(\phi, w) + P_B(\phi, w) + P_N(\phi, w) = 1$, which implies $P_T(\phi, w) + P_F(\phi, w) \le 1 - P_B(\phi, w)$.
Substituting this upper bound into the sum gives $(1 - P_B(\phi, w)) + 2P_B(\phi, w) \ge 2c_i$, which simplifies strictly to $P_B(\phi, w) \ge 2c_i - 1$.

\textbf{Tightness of the Bound (Constructibility of Epistemic Glut):}
Suppose $P_B(\phi, w) = x \ge 2c_i - 1$. Assuming the accessibility relation $R_i(w)$ is sufficiently rich to contain worlds realizing all relevant truth value assignments, I construct a valid probability distribution ensuring an epistemic glut. Let $P_T(\phi, w) = \frac{1-x}{2}$, $P_F(\phi, w) = \frac{1-x}{2}$, and $P_N(\phi, w) = 0$. Since $x \le 1$, these probabilities are non-negative and sum to $1$, making the distribution valid. 
Evaluating the positive support yields $P_T(\phi, w) + P_B(\phi, w) = \frac{1-x}{2} + x = \frac{1+x}{2}$.
Since $x \ge 2c_i - 1$, substituting for $x$ gives $\frac{1+x}{2} \ge \frac{1+2c_i-1}{2} = c_i$. 
Thus, $\mu_{i,w}(\mathcal{E}^+(\phi)) \ge c_i$. By symmetry, $\mu_{i,w}(\mathcal{E}^-(\phi)) \ge c_i$. This distribution therefore guarantees an epistemic glut.
\end{proof}

This theorem establishes a three-dimensional ``safe zone'' for epistemic agents. If the threshold $c_i = 0.9$, an agent can never hold contradictory beliefs unless the local environment itself exhibits at least a $0.8$ probability of being objectively contradictory. Notably, as $c_i \to 0.5^+$, the boundary $2c_i - 1 \to 0^+$, meaning the ``safe zone'' collapses and even minuscule ontological gluts can trigger epistemic gluts. This geometric relationship is visualized in Figure \ref{fig:paradox_boundary}.

\begin{figure}[ht]
    \centering
    \includegraphics[width=0.8\textwidth]{theorem_plot.pdf}
    \caption{The Cartography of Paradoxes: The surface represents the critical boundary $P_B = 2c_i - 1$. Points above this surface allow for Epistemic Gluts. The geometric mapping demonstrates how subjective uncertainty strictly maps onto objective bounds.}
    \label{fig:paradox_boundary}
\end{figure}

\subsection{Interpreting Ontological Glut Probabilities}
\label{sec:interpretation}

A central innovation of the 4-PEL framework is assigning a formal probability mass to the Ontological Glut state ($P_B$). A natural question arises: what does it mean for a proposition to have, for instance, a 0.6 probability of being both strictly true and strictly false simultaneously? Rather than treating $P_B$ as merely a formal algebraic artifact, it represents a precise measure of structural contradiction in the agent's evidentiary environment. Several concrete interpretations anchor this concept to real-world reasoning domains:

\begin{itemize}
    \item \textbf{Inconsistent Databases:} In federated database systems, the same primary key may map to conflicting records across different nodes. A $P_B$ of 0.6 means that 60\% of the accessible data shards concurrently assert and deny the proposition in question.
    \item \textbf{Conflicting Testimony:} In legal epistemology, independent witnesses might provide testimonies that directly contradict one another regarding a specific event. $P_B$ quantifies the probability mass of the evidence that is mutually inconsistent.
    \item \textbf{Distributed AI Systems:} When multiple specialized machine learning models (e.g., vision and language models) form an ensemble, their outputs may directly conflict. A high $P_B$ reflects an environment where the ensemble routinely produces high-confidence, contradictory classifications for the same input.
    \item \textbf{Scientific Models under Theory Conflict:} Incomplete scientific models under domain overlap, such as mutually inconsistent effective models, competing diagnostic criteria, or conflicting simulation regimes, can yield contradictory predictions. $P_B$ acts as a measure of how frequently these models overlap in producing mutually exclusive results.
    \item \textbf{Sensor Fusion:} Autonomous vehicles relying on multiple sensors (LiDAR, radar, cameras) frequently encounter conditions where one sensor asserts an obstacle is present while another asserts the path is clear. $P_B$ represents the likelihood of the sensor array entering a deadlock state.
\end{itemize}

By defining $P_B$ as the probability of encountering structurally irreconcilable evidence, the 4-PEL framework transitions from a purely mathematical object to a vital modeling tool for environments where contradiction is an objective, measurable feature of the data source.

\subsection{Why Four Values? The Necessity of Overlap}

A critical observer might ask: why use a four-valued logic? Why not use a three-valued logic, fuzzy logic, or imprecise probability to model these paradoxes?

The derivation of the Glut Boundary Theorem reveals that the result is not an arbitrary feature of the chosen framework, but emerges specifically and exclusively from the overlap structure inherent in a four-valued semantics.

\begin{theorem}[Necessity of Four-Valued Overlap]
No partition-based truth-state semantics with a standard non-negative probability measure can sustain epistemic gluts for thresholds $c > 0.5$. Any such system must either forbid simultaneous positive and negative belief, abandon finite additivity, or introduce an overlap state equivalent to $B$.
\end{theorem}
\begin{proof}
In any framework where truth states form a strict, mutually exclusive partition (such as classical bivalence or standard three-valued logic), positive support and negative support cannot overlap. If an agent experiences an epistemic glut ($P(\mathcal{E}^+) \ge c$ and $P(\mathcal{E}^-) \ge c$), and the extensions are disjoint ($\mathcal{E}^+ \cap \mathcal{E}^- = \emptyset$), then by finite additivity, the total probability measure over the union must be exactly $P(\mathcal{E}^+ \cup \mathcal{E}^-) = P(\mathcal{E}^+) + P(\mathcal{E}^-) \ge 2c$. Since $c \in (0.5, 1]$, $2c > 1$, which strictly violates the fundamental Kolmogorov axiom that the total probability mass $P(W)$ cannot exceed $1$.

One might consider alternative mathematical structures to circumvent this. For instance, \textit{signed measures} allow for negative probabilities, which could theoretically balance a sum greater than $1$. However, signed measures break the standard interpretation of subjective probability as degrees of belief. \textit{Fuzzy logic} maps propositions to a continuous truth degree $[0, 1]$, but lacks orthogonal dimensionality; a truth degree of $0.5$ implies uncertainty, not the simultaneous presence of total positive and total negative support. \textit{Probabilistic bilattices} can model intersecting supports, but they are essentially generalizations of the very 4-valued (or higher) overlap structures advocated here.

Only a 4-valued (or greater) overlap semantics natively defines a state ($B$) that satisfies both positive and negative support simultaneously ($\mathcal{E}^+ \cap \mathcal{E}^- \neq \emptyset$) within a standard, non-negative probability measure. It is precisely the mathematical size of this intersection ($P_B$) that absorbs the probability mass exceeding $1$. Thus, the $P_B \ge 2c - 1$ boundary is an exclusive consequence of structures supporting state overlap.
\end{proof}

\textbf{Computer Verification:} To ensure mathematical rigor, the core theorem and principal semantic constraints have been formalized and verified in Lean 4, utilizing rational number arithmetic and automated linear arithmetic tactics to prove the absence of edge-case anomalies.


\section{Modeling Paradoxes: Case Studies}
\label{sec:casestudies}

By embedding non-bivalent validity standards into this probabilistic framework, I can actively model and classify epistemic paradoxes rather than merely halting at contradictions, while safely blocking explosive consequences. The Lean 4 formalization explicitly organizes these paradoxes into a structural taxonomy based on their Kripke attractors:
\begin{itemize}
    \item \textbf{Class I (Self-Reference):} Paradoxes like the Liar and the G\"odel-inspired fixed points are represented as stable \textbf{Epistemic Gluts} ($B$).
    \item \textbf{Class II (Aggregation):} Paradoxes like the Lottery and Preface fall below threshold under aggregation, evaluating them as \textbf{Strict Falsity} ($F$) or Gaps.
    \item \textbf{Class III (Introspection):} Paradoxes like Moore's sentences fail self-reflection, and evaluate as \textbf{Epistemic Gaps} ($N$) under the specified model.
\end{itemize}

Table \ref{tab:paradox_resolution} summarizes the mapping of classical failures to 4-PEL classifications. The following subsections formalize these case studies, supported by Lean 4 formalizations of the core semantic and probabilistic constraints.

\begin{table}[ht]
\centering
\begin{tabular}{lll}
\hline
\textbf{Paradox} & \textbf{Classical Outcome} & \textbf{4-PEL Outcome} \\
\hline
Liar & Explosion & Stable Glut ($B$) \\
G\"odel-inspired Epistemic Fixed Points & Incompleteness (Gap) & Epistemic Glut ($B$) \\
Lottery & Closure Failure & Probabilistic Degradation ($F$) \\
Preface & Closure Failure & Epistemic Gap ($N$) \\
Moore & Incoherence & Epistemic Gap ($N$) \\
Ex Falso & Explosion & LP Blocks Inference \\
\hline
\end{tabular}
\caption{Classification of epistemic paradoxes in 4-PEL compared to classical Probabilistic Epistemic Logic.}
\label{tab:paradox_resolution}
\end{table}

\subsection{The Liar Paradox and Dynamic Fixed Points}
\begin{theorem}[ST-Validity and Dynamic Stability of the Liar]
Let $p$ be the Liar sentence and $i$ an agent at world $w$. Assuming the agent's prior assigns nonzero probability to the glut state (i.e., $P_B(p, w) > 0$), upon conditionalizing on the Liar's self-referential evidence ($E \equiv p \leftrightarrow \neg p$), the updated probability measure assigns probability 1 to the Ontological Glut state ($P_B^*(p, w) = 1$), forming a unique stable state. The resulting inference $B_i(p) \implies B_i(\neg p)$ is ST-valid.
\end{theorem}
\begin{proof}
Consider the Liar Paradox: ``This sentence is false.'' Evaluated strictly classically, this leads to an explosive contradiction because a classical measure cannot satisfy $P_T(p) = P_F(p)$ while maintaining $P_T(p)+P_F(p)=1$. To formally capture the self-referential nature of the evidence, I define the biconditional as $\phi \leftrightarrow \psi := (\phi \to \psi) \land (\psi \to \phi)$.
In 4-PEL, assigning high probability to a glut state for the Liar is not an ad-hoc stipulation. Suppose an agent begins with a mostly classical distribution but retains a small degree of prior glut probability (e.g., $P_T=0.49, P_F=0.49, P_N=0.01, P_B=0.01$). They then observe the evidence $E \equiv (p \leftrightarrow \neg p)$. Under Belnap-Dunn FDE semantics, the equivalence $p \leftrightarrow \neg p$ evaluates to positive support \textit{strictly and uniquely} in state $B$. Classical states $T$ and $F$ yield strictly false equivalences ($1 \leftrightarrow 0 = 0$). For the gap state $N = \langle 0,0 \rangle$, since $\neg N = N$, the implication $N \to N$ evaluates to $\neg N \lor N = N \lor N = N$. The conjunction $N \land N$ remains $N$, which has zero positive support.
Applying the conditionalization rule (Definition 4.1):
\begin{equation*}
\mu_{i,w}^*(\{w'\}) = \frac{\mu_{i,w}(\{w'\} \cap \mathcal{E}^+(E) \cap R_i(w))}{\mu_{i,w}(\mathcal{E}^+(E) \cap R_i(w))}
\end{equation*}
Since the intersection of $\{w'_{T, F, N}\}$ with $\mathcal{E}^+(E)$ is empty, their numerator is $0$. The agent's probability mass is entirely stripped from $T, F,$ and $N$. Since $P_B > 0$, the denominator consists solely of $P_B$, ensuring it is well-defined. Because probability mass over $A$ is conditionally restricted to the positive extension of $E$, all weight assigned to $N$ and $F$ vanishes. The sole remaining mass strictly belongs to $B$. Under this update rule, any prior with $P_B > 0$ conditionalizes onto the glut state.
Consequently, the threshold belief operator evaluates the premise $B_i(p)$ as an epistemic glut ($B$). Evaluating the inference $B_i(p) \implies B_i(\neg p)$ under the Strict-to-Tolerant (ST) validity standard reveals that the inference is vacuously \textit{valid} (since the premise is not strictly true), preventing explosion.
\end{proof}

\subsection{Epistemic Fixed Points}
\begin{theorem}[G\"odel-Inspired Epistemic Fixed Points]
Let $G$ be an epistemic fixed point defined structurally such that $G \leftrightarrow \neg B_i(G)$. If the agent $i$ demands strict certainty ($c_i = 1$) and $G$ has probability mass 1 for the ontological glut state ($B$), the 4-PEL framework absorbs $G$ into a stable Epistemic Glut ($B$).
\end{theorem}
\begin{proof}
In classical bivalent frameworks, a sentence asserting its own unprovability (analogous to G\"odel's First Incompleteness Theorem) creates a gap: neither provable nor disprovable within the system. In this epistemic analogue, let $B_i$ denote strict provability/belief ($c_i = 1$), and let $G$ be the fixed point $G \leftrightarrow \neg B_i(G)$.
If we evaluate $G$ as a strictly bivalent proposition, we arrive at a gap. However, in 4-PEL, the system can absorb the fixed point as a glut. 
I computationally verified this via Lean 4: By assigning the Ontological Glut state ($B$) to $G$ with probability 1, the positive and negative mass are both $1 \ge c_i$. Therefore, $B_i(G)$ evaluates to $B$. The negation of a glut is a glut, so $\neg B_i(G)$ evaluates to $B$. The fixed point equation holds: $v(G) = v(\neg B_i(G)) = B$. 
While this is not a formal treatment of arithmetic incompleteness, it provides a powerful case study: the example shows that a self-referential epistemic fixed point can be modeled as a stable glut rather than as a gap. A probabilistically resilient agent does not necessarily ``halt'' at self-referential unprovability; instead, it actively holds a contradictory but mathematically stable belief.
\end{proof}

\subsection{The Lottery Paradox and Deductive Closure}
\begin{example}[Failure of Deductive Closure]
There exists a 4-PEL model and Lockean threshold $c$ such that for an agent $i$, the inference $B_i(\neg p_1) \land B_i(\neg p_2) \implies B_i(\neg p_1 \land \neg p_2)$ is invalid under SS, ST, and LP validity standards.
\end{example}
The Lottery Paradox demonstrates that an agent can rationally believe of each individual ticket in a fair lottery that it will lose, yet not believe the conjunction that all tickets will lose. 
By mapping a 3-ticket lottery to the 4-PEL set of reachable worlds $R_i(w)$, I assign a uniform local probability $\mu_{i,w}$ of $1/3$ to each ticket winning. Here, the base ontology is entirely classical: each world contains exactly one winner, meaning tickets are strictly true or strictly false.
With a threshold $c=0.6$, the agent strictly believes each ticket will lose ($B_i(\neg p_1)$ and $B_i(\neg p_2)$ evaluate to strict truth $T$). However, the conjunction $\neg p_1 \land \neg p_2$ is only strictly true in one world (probability $1/3 < 0.6$). Because the probability drops below $1-c$, the aggregate belief $B_i(\neg p_1 \land \neg p_2)$ drops entirely out of $T$ and evaluates to strict falsity $F$. 
While a classical PEL can model this, 4-PEL is essential here because it provides a unified algebraic space where the aggregated conclusion falls below threshold into a non-bivalent epistemic state without breaking the formal machinery. The system explicitly flags the conjunction as probabilistically degraded, aligning perfectly with standard probabilistic epistemologies while formalizing the breakdown of aggregation within the 4-valued algebra.

\subsection{The Preface Paradox and Epistemic Gaps}
\begin{example}[Preface Paradox as Epistemic Gap]
In a model capturing the Preface conditions with Lockean threshold $c$, individual beliefs $B_i(s_k)$ evaluate to Strict Truth ($T$), while the belief in the conjunction $B_i(\bigwedge_k s_k)$ evaluates to an Epistemic Gap ($N$). The inference to the conjunction is correctly identified as invalid under ST and LP validities.
\end{example}
The Preface Paradox, much like the Lottery Paradox, illustrates a failure of deductive closure. However, unlike the independent external draws of a lottery, the Preface paradox emerges through a structurally distinct, diffuse accumulation of minor error probabilities across the agent's own authored work. An author rationally believes each individual statement $s_k$ in their meticulously researched book, yet acknowledges in the preface that the book inevitably contains at least one error.
In 4-PEL, I model this by evaluating the truth of statements across accessible worlds $R_i(w)$ where each world encapsulates exactly one error. With a Lockean threshold $c=0.85$, if the probability of the entire book being flawless is $\mu_{i,w}(\{w_{flawless}\}) = 0.7$, and individual errors split the remaining $0.3$ probability, each statement $s_k$ retains a high probability ($0.9 \ge 0.85$), making $B_i(s_k)$ evaluate to strictly true ($T$). 
However, the conjunction evaluates to positive support of $0.7$ ($< 0.85$) and negative support of $0.3$ ($< 0.85$). As a result, the belief in the conjunction $B_i(\bigwedge_k s_k)$ evaluates to an \textbf{Epistemic Gap} ($N$). The agent does not strictly disbelieve the conjunction; rather, they rationally suspend judgment. Because the conclusion is an epistemic gap, the inference from strict premises fails under ST-validity, correctly demonstrating that 4-PEL blocks the accumulation of errors into absolute truth and enforces a rational suspension of belief.

\subsection{Moore's Paradox and Introspection}
\begin{example}[Introspection of Moorean Gluts]
Let $\phi = p \land \neg B_i(p)$. If $P_B(p, w) = 1$, then $B_i(\phi)$ evaluates to an Epistemic Glut ($B$). If $P_T(p, w) = 1$, $B_i(\phi)$ evaluates to Strict Falsity ($F$).
\end{example}
Moore's Paradox involves sentences of the form: ``It is raining, but I don't believe it is raining'' ($p \land \neg B_i(p)$). A rational agent cannot coherently believe this: $B_i(p \land \neg B_i(p))$.
In 4-PEL, epistemic introspection reflects the underlying ontology. In a classical world ($P_T(p, w) = 1.0$), the sentence evaluates to strictly false ($F$), and the belief in it is also $F$, matching classical epistemology. However, if the agent inhabits an inherently contradictory world where $p$ is an ontological glut ($P_B(p, w) = 1.0$), $B_i(p)$ and its negation both evaluate as gluts. Consequently, the agent's belief in Moore's sentence also evaluates to an epistemic glut ($B$). A continuous spectrum exists between these extremes; as $P_B(p, w)$ grows, the agent smoothly transitions from classical falsity to epistemic glut. Thus, 4-PEL preserves classical boundaries in consistent worlds while allowing introspection of paradoxes in glut-heavy environments.

\subsection{Epistemic Explosion (Ex Falso Quodlibet)}
\begin{theorem}[Paraconsistency and Epistemic Explosion]
For any agent $i$, the inference $B_i(p \land \neg p) \implies B_i(q)$ is invalid under LP validity. Thus, 4-PEL effectively prevents epistemic explosion.
\end{theorem}
The principle of \textit{Ex Falso Quodlibet} dictates that believing a contradiction should lead to believing everything ($B_i(p \land \neg p) \implies B_i(q)$). A robust paraconsistent system must block this ``epistemic explosion.''
In 4-PEL, if an agent believes a contradiction (because $p$ is an ontological glut with $P_B(p, w) = 1.0$), the premise $B_i(p \land \neg p)$ evaluates to a glut ($B$). If $q$ is an unrelated, strictly false proposition ($P_F(q, w) = 1.0$), the conclusion $B_i(q)$ evaluates to strict falsity ($F$). 
While Strict-to-Tolerant (ST) validity holds vacuously (since the premise is not strictly true), the \textbf{Tolerant-to-Tolerant (LP) validity} correctly identifies this inference as \textit{invalid}. LP demands that a tolerated premise ($T$ or $B$) yields a tolerated conclusion, which $F$ is not. Therefore, LP validity actively prevents epistemic explosion, establishing 4-PEL as a genuinely paraconsistent epistemic logic.

\section{Topological Cartography and Meta-Glut Collapse}
\label{sec:topological_cartography}

The Cartography of Paradoxes (Section 4) maps the boundary where first-order epistemic gluts emerge. A crucial question arises when we consider higher-order beliefs: how deep does the paradox penetrate? By formally analyzing the geometry of nested belief operators $B_i(B_i(\ldots(p)))$, we discover profound differences between isolated agents and multi-agent topologies.

\subsection{Total Hierarchical Collapse in Isolated Agents}
Let $S_n$ denote the surface boundary for an epistemic glut at reflection depth $n$ (i.e., the state where $B_i^n(p) = B$). Intuitively, one might expect higher-order reflection to dampen contradictions, creating increasingly strict boundaries ($S_1 < S_2 < S_3 \dots$).

\begin{theorem}[The Theorem of Meta-Glut Hierarchy Collapse]
Under Euclidean and transitive accessibility ($R_i$), stable local measures, and uniform evaluation of $B_i(p)$ across $R_i(w)$, if a proposition breaches the first-order paradox boundary ($P_B \ge 2c_i - 1$) for an agent $i$ with a constant threshold $c_i$, it propagates through the entire infinitely deep meta-belief hierarchy. Formally:
$B_i(p) = B \implies \forall n \in \mathbb{N}^+, B_i^n(p) = B$.
Thus, the topological boundaries collapse perfectly onto each other: $S_1 = S_\infty$.
\end{theorem}
\begin{proof}
This phenomenon was formally verified in Lean 4. If $B_i(p)$ evaluates to $B$ at world $w$, and the agent perfectly introspects (via isolated Kripke accessibility), the probability mass on the worlds where $B_i(p)$ is a glut is $1.0$. Because $1.0 \ge c_i$, the second-order belief $B_i(B_i(p))$ also evaluates to $B$. By mathematical induction, this propagates infinitely. There are no fractal boundaries in a solitary mind; the induction propagates through all finite nesting depths simultaneously.
\end{proof}

\subsection{Multi-Agent Topological Curvature}
Topological curvature in the boundary surfaces ($S_2 \neq S_1$) is mathematically impossible to sustain in an isolated agent. It can \textit{only} be rescued by social interaction across a multi-agent system where thresholds $c_i$ vary.

\begin{theorem}[Rational Disagreement and Topological Curvature]
Let $i, j \in Ag$ be agents observing the same local probability measure ($\mu_{i,w} = \mu_{j,w}$) but holding different Lockean thresholds $c_i > c_j$. There exist models where $B_j(\phi)$ evaluates to an epistemic glut ($B$), while $B_i(\phi)$ evaluates to an epistemic gap ($N$). Agent $i$ can hold the higher-order belief $B_i(B_j(\phi))$ without holding $B_i(\phi)$.
\end{theorem}
The power of a multi-agent framework lies in modeling perspective-dependent reasoning. Suppose agents $i$ and $j$ share the same accessibility relation and observe an ontological glut probability of $P_B(\phi, w) = 0.6$, with the remaining mass purely classical: $P_T = 0.2, P_F = 0.2, P_N = 0$.
Let agent $i$ be highly risk-averse with a strict threshold $c_i = 0.9$, and agent $j$ be more permissive with $c_j = 0.6$. For agent $j$, because $0.6 \ge c_j$, their threshold belief operator evaluates $B_j(\phi)$ as an epistemic glut ($B$). However, for agent $i$, the probability falls short of $c_i$, leading $B_i(\phi)$ to evaluate as an epistemic gap ($N$). 
Because the 4-PEL framework allows seamless nesting of belief operators, agent $i$ can perfectly model agent $j$'s contradictory state. If agent $i$'s local measure places probability $1.0$ on worlds where $B_j(\phi)$ evaluates to $B$, then $B_i(B_j(\phi))$ evaluates to strict truth ($T$), even though $B_i(\phi)$ is an epistemic gap ($N$). This demonstrates how 4-PEL elegantly models rational higher-order disagreement about contradictions, rescuing the topological curvature of the meta-belief hierarchy through relativistic social epistemology.

\section{Axiomatization and Metatheory}
\label{sec:metatheorems}

To meet the rigorous standards of formal logic, I provide a Hilbert-style axiomatization for the static fragment of 4-PEL, alongside establishing soundness and detailing the reduction strategy for completeness. We first define the central non-bivalent consequence relations used throughout this work:
\begin{align*}
\Gamma \models_{ST} \phi &\iff \text{whenever every premise in } \Gamma \text{ is strictly true, } \phi \text{ is tolerantly true.} \\
\Gamma \models_{LP} \phi &\iff \text{whenever every premise in } \Gamma \text{ is tolerantly true, } \phi \text{ is tolerantly true.}
\end{align*}

\subsection{Axiomatization and Soundness}
The deductive system for 4-PEL, denoted $\vdash_{ST}$, is anchored in the Strict-Tolerant validity standard, ensuring that inferences preserve at least tolerant truth (preventing explosion). The core axioms include schemata valid in FDE, such as double negation principles, De Morgan principles, and conjunction/disjunction rules appropriate to the positive/negative support semantics. The primary inference rule is Modus Ponens ($\phi, \phi \to \psi \vdash \psi$), which is uniquely ST-valid but notably invalid under LP-validity.

To formalize this two-level proof discipline, the calculus tracks strict and tolerant statuses using typed derivability:
\begin{table}[ht]
\centering
\begin{tabular}{ll}
\hline
\textbf{Judgment} & \textbf{Meaning} \\
\hline
$\Gamma \vdash_S \varphi$ & strict derivability \\
$\Gamma \vdash_T \varphi$ & tolerant derivability \\
MP-ST & from strict $\varphi$ and strict $\varphi \to \psi$, infer tolerant $\psi$ \\
\hline
\end{tabular}
\end{table}

Representative axiom schemata governing this calculus include strict double negation ($\vdash_S \neg\neg\phi \leftrightarrow \phi$), De Morgan laws for conjunction and disjunction, and belief-distribution axioms tracking positive and negative support. Strict derivability ($\vdash_S$) operates as the generative engine for tautologies, while tolerant derivability ($\vdash_T$) acts as the absorbing boundary for conclusions. A strict premise guarantees a tolerant conclusion via ST Modus Ponens, but a tolerant premise does not propel further strict inferences, structurally neutralizing explosive derivations without stripping the logic of classical deductive power in consistent states.

\begin{theorem}[Soundness of 4-PEL]
If $\vdash_{ST} \phi$, then $\models_{ST} \phi$. Any formula provable within the Hilbert-style calculus of 4-PEL is valid under ST semantics.
\end{theorem}
\begin{proof}
The proof proceeds by structural induction on the length of the derivation. Each selected FDE-valid axiom schema in the calculus evaluates tolerantly, i.e. as $T$ or $B$, in any 4-PEL model. Crucially, because ST consequence is not standard truth preservation, ST-soundness requires a carefully typed derivability relation or a two-level proof discipline tracking strict and tolerant statuses. By ensuring that Modus Ponens preserves this tolerant evaluation property when transitioning from strictly true premises to a tolerantly true conclusion, the structural safety of the logic is maintained. This metatheorem has been fully mechanized and proven within the Lean 4 proof assistant kernel.
\end{proof}

\subsection{Decidability and Completeness via Classical Translation}
Proving strong completeness from scratch for a logic combining probability, non-bivalent evaluations, and threshold modal operators is notoriously complex. However, the static fragment of 4-PEL admits a reduction route to classical Probabilistic Epistemic Logic (PEL) defined by Fagin, Halpern, and Megiddo (1990) \cite{fagin1990}.

\begin{proposition}[Reduction of the Static Fragment]
This establishes a conditional route to decidability and strong completeness for the static fragment, pending a fully explicit compositional proof of the translation.
\end{proposition}
\begin{proof}[Proof Sketch]
For any atomic proposition $p \in \Phi$, we introduce two strictly classical (bivalent) atomic propositions: $p^+$ (representing positive support) and $p^-$ (representing negative support). The 4-valued space $W$ is mathematically isomorphic to a classical probability space constructed over the powerset of $\{p^+, p^- : p \in \Phi\}$. 
We define recursive translations $\tau^+(\phi)$ and $\tau^-(\phi)$ such that:
\begin{align*}
\mathcal{M}, w \models^+ \phi &\iff \mathcal{M}^\tau, w \models \tau^+(\phi) \\
\mathcal{M}, w \models^- \phi &\iff \mathcal{M}^\tau, w \models \tau^-(\phi)
\end{align*}

Crucially, the threshold belief operator $B_i(\phi)$ translates directly to a classical probability inequality via this recursion: $\tau^+(B_i(\phi)) := P(\tau^+(\phi)) \ge c$. Because the translated formulas are strictly classical PEL formulas constructed over linear weight inequalities, 4-PEL inherits the standard Fagin-Halpern canonical model construction. Thus, any 4-PEL formula valid in all models is provable by translating it to classical PEL, applying the complete classical PEL axioms, and translating back. This would yield decidability and a reduction pathway to strong completeness once the compositional translation is fully established. The direct, axiomatic completeness of the dynamic conditionalization fragment remains an open area for future work.
\end{proof}

\subsection{Complexity of Model Checking}
It is crucial to establish the computational properties of the underlying system. The 4-PEL models as defined are finite Kripke structures equipped with finite sets of agents and propositions. 
\begin{theorem}[Model Checking Complexity]
Given a finite 4-PEL model $\mathcal{M}$ and a formula $\phi$, determining whether $\mathcal{M}, w \models^+ \phi$ or $\mathcal{M}, w \models^- \phi$ operates in polynomial time $\mathcal{O}((|W| + |R|) \times |\phi|)$, where $|W|$ is the number of worlds, $|R|$ represents the accessibility transitions, and $|\phi|$ is the length of the formula, assuming unit-cost arithmetic for the rational probability threshold evaluations.
\end{theorem}
\begin{proof}
Model checking is computationally tractable. The evaluation of standard FDE truth functions at any single world takes $\mathcal{O}(|\phi|)$ steps. The threshold belief operator $B_i(\phi)$ requires traversing the accessibility relation $R_i$ to evaluate $\phi$ across reachable worlds and summing their pre-computed local probability masses $\mu_{i,w}$. Employing bottom-up dynamic programming (memoization) ensures that truth values for subformulas are computed exactly once per world, yielding the standard modal model-checking bound of $\mathcal{O}((|W| + |R|) \times |\phi|)$. This establishes 4-PEL as highly efficient for computational applications like AI belief aggregation.
\end{proof}

The Lean 4 formalization verifies the stated core probabilistic constraints and principal validity inferences under the definitions used here.

\subsection{Scope of Formal Verification}

While this paper heavily utilizes the Lean 4 interactive theorem prover, it is critical to precisely delineate what has and has not been formally verified. This prevents the conflation of mathematical proof with philosophical assertion.

\textbf{Lean 4 formally verifies:}
\begin{itemize}
    \item \textbf{Theorem Derivations:} The mathematical validity of the Glut Boundary Theorem ($P_B \ge 2c - 1$) given the stated definitions.
    \item \textbf{Semantic Consistency:} The structural absence of explosive contradictions when applying ST and LP validity standards to the 4-valued threshold belief operator.
    \item \textbf{Probability Inequalities:} The preservation of Kolmogorov probability axioms across the generalized conditionalization rule.
\end{itemize}

To provide exact traceability to the formal proofs in the repository, Table \ref{tab:lean_verification} maps the theoretical claims in this manuscript to their explicit Lean 4 theorem names and the structural assumptions required for the kernel to accept the proofs.

\begin{table}[ht]
\centering
\resizebox{\textwidth}{!}{%
\begin{tabular}{lllc}
\hline
\textbf{Paper Theorem} & \textbf{Lean Theorem Name} & \textbf{Structural Assumptions} & \textbf{Verified?} \\
\hline
Glut Boundary & \texttt{glut\_boundary} & Finite measure, Threshold $c > 1/2$ & Yes \\
Tightness & \texttt{glut\_boundary\_tight} & Rich state space yielding valid distribution & Yes \\
ST Soundness & \texttt{st\_soundness\_valid} & ST-typed derivability calculus & Yes \\
Conditionalization & \texttt{cond\_measure\_valid} & Denominator non-zero ($P(E^+) > 0$) & Yes \\
\hline
\end{tabular}%
}
\caption{Mapping of manuscript theorems to exact Lean 4 formalizations.}
\label{tab:lean_verification}
\end{table}
\textbf{Lean 4 does not verify:}
\begin{itemize}
    \item \textbf{Philosophical Interpretations:} The mappings of Ontological Gluts to real-world phenomena (e.g., inconsistent databases or sensor fusion).
    \item \textbf{Empirical Applicability:} Whether actual human agents or AI systems naturally conform to Lockean threshold dynamics.
    \item \textbf{Choice of Primitives:} The philosophical justification for selecting the Belnap-Dunn 4-valued ontology over other non-classical structures.
\end{itemize}

By strictly separating the verifiable mathematical architecture from its interpretative applications, the 4-PEL framework ensures rigorous logical foundations.


\section{Conclusion}
\label{sec:conclusion}

This paper introduces the 4-Valued Probabilistic Epistemic Logic (4-PEL), a framework that unifies the quantitative nuance of Probabilistic Epistemic Logic with the qualitative resilience of Belnap-Dunn 4-valued semantics. By formally charting the ``Cartography of Paradoxes,'' I demonstrated mathematically and computationally (via Lean 4) that subjective epistemic contradictions are safely bounded by objective ontological gluts ($P_B \ge 2c - 1$). The formal verification extends beyond specific models to general theorems, proving that Strict-to-Tolerant (ST) validity preserves classical truth in consistent states, while Tolerant-to-Tolerant (LP) validity blocks epistemic explosion (ex falso quodlibet).

Furthermore, by applying these non-bivalent validities to probabilistic belief operators, 4-PEL provides a robust mechanism to model and classify five paradoxical case studies, while blocking explosive consequences in the relevant paraconsistent settings. The framework also demonstrates multi-agent rational disagreement. The seamless transition of classical contradictions into rational epistemic gaps or bounded gluts showcases the framework's expressive power. 

\section{Future Work}
\label{sec:futurework}

Future research will explore extending 4-PEL into full Dynamic Epistemic Logic (DEL) \cite{baltag2008probabilistic} with relational updates. Additionally, the polynomial-time model checking complexity suggests potential applications in artificial intelligence. As AI agents aggregate crowdsourced datasets, they encounter inherently contradictory training data. As an interpretive application, implementing 4-PEL as a foundational layer could allow artificial systems to safely compartmentalize contradictions and maintain paraconsistent reasoning under uncertainty.
\section{Data Availability}
The Lean 4 code formalizing the 4-PEL framework, the Glut Boundary Theorem, and the core semantic/probabilistic constraints underlying the paradox case studies is included in the supplementary material provided with this submission and accessible via the public repository: \url{https://github.com/cr4bbz/4PEL-Lean-Formalization}.

\appendix
\section{Lean 4 Formalization Snippets}
To demonstrate the rigor of the computational verification, the following snippets highlight the core formalization of the 4-PEL framework in Lean 4.

\subsection{4-Valued Truth and the 4-PEL Model}
I formalize the Belnap-Dunn logic using a boolean pair representing positive and negative support.
\begin{lstlisting}[language=Caml]
structure FDEValue where
  pos : Bool
  neg : Bool

def T : FDEValue := { pos := true, neg := false }
def B : FDEValue := { pos := true, neg := true }
def N : FDEValue := { pos := false, neg := false }
def F : FDEValue := { pos := false, neg := true }

structure Model (W Ag Atom : Type) where
  R : Ag -> W -> Set W
  mu : Ag -> W -> Set W -> Rat
  val : W -> Atom -> FDEValue
  c : Ag -> Rat
  -- Axioms ensuring mu is a valid probability measure over R(w) omitted for brevity.
  -- The full formalization of non-negativity, additivity, and total mass = 1 
  -- can be found in the Model.lean file in the public repository.
\end{lstlisting}

\subsection{The Threshold Belief Operator and Soundness}
The belief operator integrates the local measure with the evaluated positive/negative extensions of any abstract formula.
\begin{lstlisting}[language=Caml]
def belief (m : Model W Ag Atom) (i : Ag) (w : W) (v_phi : W -> FDEValue) : FDEValue :=
  let prob_pos := m.mu i w {w' \in m.R i w | (v_phi w').pos = true}
  let prob_neg := m.mu i w {w' \in m.R i w | (v_phi w').neg = true}
  { pos := prob_pos >= m.c i, neg := prob_neg >= m.c i }
\end{lstlisting}
Beyond these semantic evaluations, the repository now features a complete structural axiomatization of the logic. The syntax is formalized as an inductive datatype (\texttt{Formula}) alongside syntactic macros for disjunction and material implication, and we provide a safe Hilbert-style calculus (\texttt{ST\_Entails}). Crucially, the metatheoretical ST-Soundness theorem (\(\vdash_{ST} \phi \implies \models_{ST} \phi\)) is fully verified by the Lean 4 kernel via structural induction, guaranteeing that all provable inferences safely preserve tolerant truth within our 4-valued threshold semantics.

\begin{thebibliography}{99}

\bibitem{pan2024probabilistic}
Yixin Pan, Meiyun Guo.
\textit{Probabilistic epistemic logic based on neighborhood semantics}.
Synthese (2024) 203:135.

\bibitem{baltag2008probabilistic}
Alexandru Baltag, Sonja Smets.
\textit{A Qualitative Theory of Dynamic Interactive Belief Revision}.
Texts in Logic and Games 3, Amsterdam University Press (2008).

\bibitem{vanbenthem2011}
Johan van Benthem.
\textit{Logical Dynamics of Information and Interaction}.
Cambridge University Press (2011).

\bibitem{shafer1976}
Glenn Shafer.
\textit{A Mathematical Theory of Evidence}.
Princeton University Press (1976).

\bibitem{list2002}
Christian List, Philip Pettit.
\textit{Aggregating Sets of Judgments: An Impossibility Result}.
Economics and Philosophy, 18(1):89--110 (2002).

\bibitem{barrio2025non}
Eduardo Barrio, Camillo Fiore, Federico Pailos.
\textit{Non-Bivalent Validity}.
Studia Logica (2025).

\bibitem{belnap1977}
Nuel D. Belnap.
\textit{A Useful Four-Valued Logic}.
In: J.M. Dunn, G. Epstein (eds) Modern Uses of Multiple-Valued Logic. Springer, Dordrecht (1977).

\bibitem{dunn1976}
J. Michael Dunn.
\textit{Intuitive Semantics for First-Degree Entailments and 'Coupled Trees'}.
Philosophical Studies, 29(3):149--168 (1976).

\bibitem{cobreros2012}
Pablo Cobreros, Paul Egré, David Ripley, Robert van Rooij.
\textit{Tolerant, Classical, Strict}.
Journal of Philosophical Logic, 41(2):347--385 (2012).

\bibitem{priest1979}
Graham Priest.
\textit{The Logic of Paradox}.
Journal of Philosophical Logic, 8(1):219--241 (1979).

\bibitem{kyburg1961}
Henry E. Kyburg Jr.
\textit{Probability and the Logic of Rational Belief}.
Wesleyan University Press (1961).

\bibitem{makinson1965}
David C. Makinson.
\textit{The Paradox of the Preface}.
Analysis, 25(6):205--207 (1965).

\bibitem{fagin1995}
Ronald Fagin, Joseph Y. Halpern, Yoram Moses, Moshe Y. Vardi.
\textit{Reasoning About Knowledge}.
MIT Press (1995).

\bibitem{fagin1990}
Ronald Fagin, Joseph Y. Halpern, Nimrod Megiddo.
\textit{A Logic for Reasoning About Probabilities}.
Information and Computation, 87(1-2):78--128 (1990).

\bibitem{hintikka1962}
Jaakko Hintikka.
\textit{Knowledge and Belief: An Introduction to the Logic of the Two Notions}.
Cornell University Press (1962).

\bibitem{levi1980}
Isaac Levi.
\textit{The Enterprise of Knowledge}.
MIT Press (1980).

\bibitem{hawthorne2004}
John Hawthorne.
\textit{Knowledge and Lotteries}.
Oxford University Press (2004).

\bibitem{douven2021}
Igor Douven.
\textit{The Epistemology of Indicative Conditionals}.
Cambridge University Press (2021).

\bibitem{priest2006}
Graham Priest.
\textit{Doubt Truth to be a Liar}.
Oxford University Press (2006).

\bibitem{jager2001}
Gerhard J\"ager.
\textit{Some Notes on Probabilistic Modal Logic}.
In: Logic, Language and Computation, Springer (2001).

\end{thebibliography}

\end{document}
